% !TEX program = xelatex
% -*- coding: utf-8 -*-
\documentclass[11pt,a4paper]{article}

% ===== 패키지 =====
\usepackage{fontspec}
\usepackage{amsmath, amssymb, amsthm}
\usepackage{geometry}
\usepackage{hyperref}
\usepackage{enumitem}
\usepackage{booktabs}
\usepackage{array}
\usepackage{xcolor}
\usepackage[most]{tcolorbox}
\usepackage{titlesec}
\usepackage{fancyhdr}

% ===== 폰트 (Noto Sans/Serif CJK KR) =====
% PostScript 이름으로 ttc 안의 특정 폰트 지정
\setmainfont{NotoSerifCJKkr-Regular}[
  BoldFont = NotoSerifCJKkr-Bold,
  ItalicFont = NotoSerifCJKkr-Regular,
  ItalicFeatures = {FakeSlant=0.2},
  BoldItalicFont = NotoSerifCJKkr-Bold,
  BoldItalicFeatures = {FakeSlant=0.2}
]
\setsansfont{NotoSansCJKkr-Regular}[
  BoldFont = NotoSansCJKkr-Bold,
  ItalicFont = NotoSansCJKkr-Regular,
  ItalicFeatures = {FakeSlant=0.2}
]

% ===== 페이지 설정 =====
\geometry{
  a4paper,
  left=25mm, right=25mm, top=28mm, bottom=28mm,
  headheight=14pt
}
\linespread{1.35}

% ===== 색상 =====
\definecolor{accent}{HTML}{378ADD}
\definecolor{accentdark}{HTML}{185FA5}
\definecolor{boxbg}{HTML}{F1EFE8}
\definecolor{boxborder}{HTML}{B4B2A9}
\definecolor{quotebg}{HTML}{FAEEDA}
\definecolor{quoteborder}{HTML}{BA7517}

% ===== 하이퍼링크 =====
\hypersetup{
  colorlinks=true,
  linkcolor=accentdark,
  urlcolor=accentdark,
  citecolor=accentdark,
  pdftitle={달랑베르에서 PINN까지: 편미분방정식의 270년 여정},
  pdfauthor={Claude \& Crowdy}
}

% ===== 섹션 스타일 =====
\titleformat{\section}
  {\Large\bfseries\color{accentdark}}
  {\thesection.}{0.6em}{}
\titleformat{\subsection}
  {\large\bfseries\color{accentdark}}
  {\thesubsection.}{0.5em}{}
\titleformat{\subsubsection}
  {\normalsize\bfseries}
  {\thesubsubsection.}{0.4em}{}

% ===== 헤더/푸터 =====
\pagestyle{fancy}
\fancyhf{}
\fancyhead[L]{\small\itshape 달랑베르에서 PINN까지}
\fancyhead[R]{\small\thepage}
\renewcommand{\headrulewidth}{0.4pt}

% ===== 인용 박스 =====
\newtcolorbox{quotebox}{
  colback=quotebg,
  colframe=quoteborder,
  boxrule=0.5pt,
  arc=2pt,
  left=10pt, right=10pt, top=8pt, bottom=8pt,
  before skip=10pt, after skip=10pt
}

\newtcolorbox{notebox}{
  colback=boxbg,
  colframe=boxborder,
  boxrule=0.5pt,
  arc=2pt,
  left=10pt, right=10pt, top=8pt, bottom=8pt,
  before skip=8pt, after skip=8pt
}

% ===== 문서 시작 =====
\begin{document}

\begin{titlepage}
\centering
\vspace*{3cm}

{\Huge\bfseries 달랑베르에서 PINN까지\par}
\vspace{0.5cm}
{\LARGE\itshape 편미분방정식의 270년 여정\par}
\vspace{2cm}

{\large\color{accentdark}
\rule{0.7\textwidth}{0.5pt}\\
\vspace{0.5cm}
진동하는 줄 한 가닥에서 시작된 이야기가\\
오늘날 환자의 심장 안 혈류 시뮬레이션과\\
신경망 기반 수치해석으로 이어지기까지\\
\vspace{0.5cm}
\rule{0.7\textwidth}{0.5pt}\par
}

\vfill

{\large Crowdy \& Claude 대화록\par}
\vspace{0.5cm}
{\normalsize 2026년 5월\par}

\end{titlepage}

\tableofcontents
\newpage

% ====================================================================
\section{서문}
% ====================================================================

이 문서는 18세기 프랑스의 수학자 장 르 롱 달랑베르(Jean le Rond d'Alembert, 1717--1783)에서 시작하여, 그가 던진 진동하는 줄의 문제가 어떻게 푸리에 급수, 나비에-스토크스 방정식, 그리고 현대의 물리 정보 신경망(PINN)으로 이어졌는지를 따라가는 대화의 기록이다.

본 대화의 핵심 주제는 다음과 같다.

\begin{itemize}[leftmargin=*, itemsep=2pt]
  \item 달랑베르라는 인물과 그의 원리
  \item 파동방정식과 그 해법
  \item 편미분방정식의 의미와 푸리에 급수 논쟁
  \item 나비에-스토크스 방정식과 심혈관 시뮬레이션
  \item 유한체적법(FVM)과 물리 정보 신경망(PINN)의 비교
\end{itemize}

% ====================================================================
\section{장 르 롱 달랑베르라는 인물}
% ====================================================================

\subsection{어떤 사람이었는가}

18세기 프랑스 계몽주의를 대표하는 수학자, 물리학자, 철학자이다. 흔히 드니 디드로(Diderot)와 함께 \emph{백과전서(Encyclopédie)}를 공동 편집한 인물로 기억되며, 고전역학에서는 \textbf{달랑베르 원리(d'Alembert's principle)}로 그 이름이 남아 있다.

라그랑주가 해석역학을 세우기 전에, 뉴턴의 벡터적 운동방정식을 ``에너지·변분'' 관점으로 옮겨가는 결정적 다리를 놓은 사람이라고 볼 수 있다.

\subsection{``구속력은 가상 변위에 대해 일을 하지 않는다'': 달랑베르 원리}

이 문장은 달랑베르 원리의 \textbf{근본 가정}이자, 라그랑주 역학 전체가 작동하는 출발점이다.

\textbf{구속력(constraint force)}은 계의 자유로운 운동을 제한하기 위해 작용하는 힘이다. 예를 들어:
\begin{itemize}[leftmargin=*, itemsep=1pt]
  \item 책상 위 책에 작용하는 수직항력
  \item 진자 줄의 장력
  \item 철사에 꿰어진 구슬에 철사가 가하는 반력
  \item 강체 내부의 결합력
\end{itemize}

\textbf{가상 변위(virtual displacement, $\delta \mathbf{r}$)}는 어떤 한순간에 구속조건과 모순되지 않는 범위에서 가상적으로 허용된 미소 변위를 말한다. 실제 시간 흐름과 무관한, 사고 실험상의 변위이다.

핵심은 이것이다 --- 구속력은 거의 항상 \textbf{허용된 운동 방향에 수직}으로 작용한다.

\begin{itemize}[leftmargin=*, itemsep=1pt]
  \item 진자 줄의 장력은 줄 방향이고, 진자추가 움직일 수 있는 방향은 줄에 수직인 호 방향 $\Rightarrow$ 내적이 0
  \item 책상의 수직항력은 위쪽이고, 책이 미끄러질 수 있는 방향은 수평 $\Rightarrow$ 내적이 0
\end{itemize}

즉 일 $W = \mathbf{F}_{\text{구속}} \cdot \delta \mathbf{r} = 0$. 이를 통틀어 \textbf{이상적 구속(ideal constraint)}이라 부른다.

이 가정 덕분에 운동방정식
\[
\sum_i (\mathbf{F}_i^{\text{외력}} + \mathbf{F}_i^{\text{구속}} - m_i \mathbf{a}_i) \cdot \delta \mathbf{r}_i = 0
\]
에서 구속력 항이 통째로 사라진다. 결과적으로 \textbf{구속력을 미리 알지 못해도} 운동방정식을 세울 수 있게 되며, 이 우회로가 곧 라그랑주 방정식과 해밀턴 역학으로 이어진다.

\subsection{그의 일생}

\textbf{버려진 아이로 시작했다.} 1717년 11월 16일 파리의 생장르롱(Saint-Jean-le-Rond) 성당 계단에 한 갓난아기가 버려져 있었다. 이름이 없어 성당 이름을 따서 ``장 르 롱''이라 불렸다. 생모는 살롱 주인으로 유명했던 탕생 후작부인(Madame de Tencin), 생부는 포병 장교 데투슈(Louis-Camus Destouches)였다. 생모는 평생 그를 외면했지만, 생부는 익명으로 양육비와 교육비를 댔다.

\textbf{학문적 경로는 빨랐다.} 1741년 23세에 파리 과학아카데미 회원이 되었고, 1743년 26세에 발표한 \emph{동역학론(Traité de dynamique)}에서 오늘날 그의 이름이 붙은 원리를 제시했다. 이후 유체역학(달랑베르의 역설), 편미분방정식(파동방정식의 해), 천체역학(춘분점의 세차운동 설명) 등에 굵직한 기여를 남긴다.

\textbf{계몽주의의 얼굴이 되었다.} 1751년부터 디드로와 함께 \emph{백과전서}를 편집했고, 그 유명한 ``서설(Discours préliminaire)''을 직접 썼다. 1772년 아카데미 프랑세즈의 종신서기가 된다. 프로이센의 프리드리히 2세와 러시아의 예카테리나 2세가 거듭 자국으로 모셔가려 했지만 모두 거절하고 파리에 머물렀다.

\textbf{개인적 삶은 외로웠다.} 인생의 가장 깊은 애정은 살롱의 안주인 쥘리 드 레스피나스(Julie de Lespinasse)에게로 향했지만, 그녀의 마음은 다른 사람들에게 가 있었다. 1776년 그녀가 죽은 뒤 달랑베르는 큰 충격에 빠졌고, 1783년 10월 29일 파리에서 사망했다. 교회의 매장을 거부당해 표지 없는 무덤에 묻혔다 --- 성당 계단에서 시작해 무명의 무덤에서 끝난 셈이다.

% ====================================================================
\section{달랑베르의 역설과 파동방정식}
% ====================================================================

\subsection{달랑베르의 역설}

1752년 달랑베르는 \emph{유체의 저항에 관한 새 이론(Essai d'une nouvelle théorie de la résistance des fluides)}에서 다음 결과를 증명했다.

\begin{quotebox}
\textbf{이상유체(비점성·비압축·정상·비회전 흐름) 속에 놓인 강체는 그것이 어떤 모양이든 항력(drag)을 받지 않는다.}
\end{quotebox}

즉 균일한 흐름 속의 공이든 원기둥이든 비행기 모양이든, 이론상으로는 흐름이 가하는 알짜 힘이 0이라는 결론이다. 이것이 ``역설''이라 불린 이유는 명백하다 --- 강에 손을 담그면 분명히 힘을 받기 때문이다.

이상유체 가정 아래 원기둥 주위의 흐름을 풀어보면, 유선이 \textbf{앞뒤로 완전히 대칭}으로 흐른다. 베르누이 방정식에 따라 압력은 속도가 느린 곳(정체점)에서 높고, 빠른 곳에서 낮은데, 이 압력 분포 역시 앞뒤로 똑같이 대칭이다. 압력을 원기둥 표면 전체에 대해 적분하면 앞면과 뒷면이 정확히 상쇄된다.

이 역설은 단순한 계산 오류가 아니라 \textbf{이상유체 모델 자체의 결함}을 드러낸 사건이었다. 해결의 실마리는 \textbf{점성(viscosity)}, \textbf{경계층(boundary layer)}, \textbf{후류(wake)}의 세 가지로 모인다. 1904년 프란틀(Ludwig Prandtl)이 결정적 통찰을 내놓는다 --- 점성이 아무리 작아도 물체 표면 바로 옆의 얇은 층에서는 점성이 지배적이며, 이 층이 물체 뒤쪽에서 흐름으로부터 \textbf{분리(separation)}된다는 것이다.

\subsection{1차원 파동방정식}

1747년 달랑베르는 양 끝이 고정된 진동하는 줄(예: 바이올린 현)의 운동을 분석하면서 다음 편미분방정식에 도달했다.

\begin{equation}
\frac{\partial^2 u}{\partial t^2} = c^2 \frac{\partial^2 u}{\partial x^2}
\label{eq:wave1d}
\end{equation}

여기서 $u(x,t)$는 위치 $x$, 시간 $t$에서 줄의 변위, $c$는 파동의 전파 속도이다. 이것이 \textbf{1차원 파동방정식}이며, 수리물리학에서 편미분방정식을 본격적으로 사용한 거의 최초의 사례로 꼽힌다.

\subsection{달랑베르의 해법}

달랑베르의 천재성은 풀이 방법에 있었다. 그는 다음 새 좌표를 도입했다.
\[
\xi = x - ct, \quad \eta = x + ct
\]

연쇄법칙을 적용하면 파동방정식이 놀랍도록 간단한 형태로 변한다.
\[
\frac{\partial^2 u}{\partial \xi\, \partial \eta} = 0
\]

이를 적분하면 일반해가 단번에 나온다.

\begin{quotebox}
\centering
\textbf{달랑베르의 해}
\[
u(x, t) = f(x - ct) + g(x + ct)
\]
\end{quotebox}

물리적 의미는 직관적이다.
\begin{itemize}[leftmargin=*, itemsep=1pt]
  \item $f(x - ct)$: 형태를 유지한 채 \textbf{오른쪽으로 속도 $c$로 이동}하는 파동
  \item $g(x + ct)$: 형태를 유지한 채 \textbf{왼쪽으로 속도 $c$로 이동}하는 파동
\end{itemize}

즉, \textbf{모든 1차원 파동은 좌우로 진행하는 두 파동의 중첩}으로 표현된다.

초기조건 $u(x,0) = \phi(x)$와 $u_t(x,0) = \psi(x)$가 주어진 무한히 긴 줄의 경우 명시적 공식은 다음과 같다.
\begin{equation}
u(x, t) = \tfrac{1}{2}\bigl[\phi(x - ct) + \phi(x + ct)\bigr] + \frac{1}{2c}\int_{x - ct}^{x + ct} \psi(s)\, ds
\end{equation}

% ====================================================================
\section{편미분방정식이란 무엇인가}
% ====================================================================

\subsection{편미분의 의미}

고등학교 미분 $\dfrac{dy}{dx}$는 변수가 $x$ 하나뿐이다. 그런데 진동하는 줄의 변위 $u$는 두 가지에 따라 달라진다 --- \textbf{위치 $x$}와 \textbf{시간 $t$}.

이 경우 ``변화율''도 두 종류가 필요하다.

\begin{itemize}[leftmargin=*, itemsep=2pt]
  \item $\dfrac{\partial u}{\partial t}$: 한 위치를 고정하고 시간만 흐르게 둘 때 변위의 변화율 = 그 점의 \textbf{속도}
  \item $\dfrac{\partial u}{\partial x}$: 시간을 정지시키고 줄을 따라 옆으로 이동할 때 변위의 변화율 = 그 순간 줄 모양의 \textbf{기울기}
\end{itemize}

$d$ 대신 $\partial$(라운드 디)를 쓰는 이유는 \textbf{``다른 변수는 잠시 상수로 고정하고, 하나만 변화시킨 미분''}임을 표시하기 위해서다. 이를 \textbf{편미분(partial derivative)}이라고 부른다.

\subsection{파동방정식의 물리적 의미}

식 \eqref{eq:wave1d}이 말하는 바는 다음과 같다.

\begin{quotebox}
\textbf{줄의 한 점이 받는 가속도는, 그 점에서 줄이 휘어진 정도에 비례한다.}
\end{quotebox}

직관적으로, 줄을 살짝 잡아당겨 놓으면 그 부분만 볼록하게 휘어 있고, 양옆의 줄이 \textbf{원래대로 펴려고 잡아당긴다}(장력 때문). 많이 휘어 있을수록 더 세게 잡아당기고, 더 큰 가속도를 만든다.

사실상 뉴턴의 $F = ma$를 줄의 모든 점에 동시에 적용한 식이다.

\subsection{변수가 많아지는 예}

편미분방정식의 변수 개수는 ``그 현상을 위치 지정하는 데 필요한 좌표의 개수''에 해당한다.

\begin{table}[h]
\centering
\small
\begin{tabular}{lcl}
\toprule
\textbf{현상} & \textbf{변수 개수} & \textbf{변수} \\
\midrule
진동하는 줄 (1차원 파동)      & 2  & $x, t$ \\
가는 막대의 열 전도         & 2  & $x, t$ \\
진동하는 막 (북)         & 3  & $x, y, t$ \\
3차원 열 전도, 전자기파      & 4  & $x, y, z, t$ \\
유체역학 (나비에-스토크스)     & 4  & $x, y, z, t$ \\
두 입자 양자역학         & 7  & $x_1,y_1,z_1, x_2,y_2,z_2, t$ \\
볼츠만 방정식           & 7  & 위치 3 + 속도 3 + 시간 1 \\
단백질의 양자 기술        & 수만 개 & 모든 원자의 좌표 + 시간 \\
\bottomrule
\end{tabular}
\caption{여러 물리 현상의 편미분방정식 변수 수}
\end{table}

3차원 열방정식은 다음과 같다.
\[
\frac{\partial T}{\partial t} = \alpha \left( \frac{\partial^2 T}{\partial x^2} + \frac{\partial^2 T}{\partial y^2} + \frac{\partial^2 T}{\partial z^2} \right)
\]

\textbf{차원의 저주(curse of dimensionality):} 변수가 늘 때마다 메쉬 기반 계산량이 기하급수적으로 늘어나기 때문에, 다입자 양자역학이나 단백질 시뮬레이션에서는 정확한 해가 사실상 불가능하고 근사법이나 신경망을 동원한다.

% ====================================================================
\section{편미분방정식의 황금기와 난장판}
% ====================================================================

달랑베르 이후 약 150년은 \textbf{편미분방정식 풀이의 대난장판}이라 할 만했다. 그러나 이 난장판이 현대 수학과 물리학을 만들었다.

\subsection{풀이를 둘러싼 18세기의 논쟁}

\textbf{달랑베르(1747)}: 매끄러운 함수만 해로 인정.

\textbf{오일러(1748)}: 손가락으로 줄을 꺾어 튕긴 V자 모양도 해여야 한다고 반박. 함수 개념 자체의 확장을 도발.

\textbf{다니엘 베르누이(1753)}: 완전히 다른 해를 들고 등장한다.
\begin{equation}
u(x, t) = \sum_{n=1}^{\infty} A_n \sin\!\left(\frac{n\pi x}{L}\right) \cos\!\left(\frac{n\pi c t}{L}\right)
\end{equation}
모든 진동이 기본 진동과 그 정수배 배음들의 무한 합으로 표현된다는 발상.

\textbf{라그랑주(1759)}: 줄을 $N$개의 작은 입자가 연결된 것으로 보고 풀어, $N \to \infty$ 극한에서 베르누이의 해에 거의 도달하지만 마지막 단계에서 멈춘다.

\subsection{푸리에의 폭탄 (1807, 1822)}

장 바티스트 푸리에가 열방정식을 풀면서 무자비하게 선언한다.

\begin{quotebox}
\textbf{``임의의 함수는 사인과 코사인의 무한 합으로 표현된다. 이상.''}
\end{quotebox}

1807년 파리 과학아카데미에 제출된 그의 논문은 라그랑주를 포함한 심사위원들에게 \textbf{거부당한다}. 흥미로운 사실 --- 라그랑주는 푸리에가 어렸을 때 군사학교 교관이었다. 자기가 가르친 학생의 논문을 30년 뒤 심사위원으로 거부한 셈이다.

\subsection{푸리에 급수의 원리: 직교성}

3차원 공간에서 임의의 벡터를 x, y, z축 성분으로 분해할 수 있는 것처럼, 사인과 코사인 함수들 $\{1, \sin x, \cos x, \sin 2x, \cos 2x, \ldots\}$이 서로 \textbf{직교하는 무한 차원의 좌표축} 역할을 한다.

함수의 ``내적''은 곱해서 적분하는 것으로 정의되며, 한 주기에 걸쳐 적분해보면
\[
\int_0^{2\pi} \sin(mx) \sin(nx)\, dx = 0 \quad (m \neq n)
\]

임의 함수 $f(x)$의 분해:
\begin{equation}
f(x) = \frac{a_0}{2} + \sum_{n=1}^{\infty} \bigl[ a_n \cos(nx) + b_n \sin(nx) \bigr]
\end{equation}

계수는 다음 공식으로 얻는다.
\[
a_n = \frac{1}{\pi}\int_0^{2\pi} f(x) \cos(nx)\, dx, \qquad b_n = \frac{1}{\pi}\int_0^{2\pi} f(x) \sin(nx)\, dx
\]

\subsection{사각파의 예}

뾰족한 사각파를 매끄러운 사인 함수들의 합으로 표현하면:
\[
\text{사각파}(x) = \frac{4}{\pi}\left[ \sin x + \frac{1}{3}\sin 3x + \frac{1}{5}\sin 5x + \frac{1}{7}\sin 7x + \cdots \right]
\]

항을 늘릴수록 사각파에 가까워지지만, 불연속점 근처에 약 9\%의 오버슈트가 끝까지 남는다. 이를 \textbf{깁스 현상(Gibbs phenomenon)}이라 한다.

\subsection{19세기의 정리 작업}

푸리에가 옳다는 것을 엄밀하게 증명하기 위해 함수, 극한, 적분의 정의를 다시 세워야 했다.

\begin{itemize}[leftmargin=*, itemsep=1pt]
  \item \textbf{코시 (1820년대)}: 극한과 연속성의 엄밀한 정의
  \item \textbf{디리클레 (1829)}: 푸리에 급수의 수렴 조건, 함수의 현대적 정의
  \item \textbf{리만 (1854)}: 리만 적분
  \item \textbf{르베그 (1902)}: 르베그 적분, $L^2$ 공간
  \item \textbf{소볼레프 (1930년대)}: 약한 도함수
  \item \textbf{슈바르츠 (1945)}: 초함수(distribution), 디랙 델타의 엄밀한 미분
\end{itemize}

% ====================================================================
\section{나비에-스토크스 방정식}
% ====================================================================

\subsection{유체의 뉴턴 방정식}

비압축성 유체의 나비에-스토크스 방정식:
\begin{equation}
\rho\left(\frac{\partial \vec{v}}{\partial t} + (\vec{v} \cdot \nabla)\vec{v}\right) = -\nabla p + \mu \nabla^2 \vec{v} + \vec{f}
\label{eq:ns}
\end{equation}

질량 보존 조건:
\begin{equation}
\nabla \cdot \vec{v} = 0
\end{equation}

식의 구조는 \textbf{질량 $\times$ 가속도 $=$ 힘들의 합}이다.

\textbf{좌변: $ma$}
\begin{itemize}[leftmargin=*, itemsep=1pt]
  \item $\rho$: 유체 밀도
  \item $\dfrac{\partial \vec{v}}{\partial t}$: 지역 가속도
  \item $(\vec{v} \cdot \nabla)\vec{v}$: 대류 가속도 (비선형 항)
\end{itemize}

\textbf{우변: 힘들}
\begin{itemize}[leftmargin=*, itemsep=1pt]
  \item $-\nabla p$: 압력차에 의한 힘
  \item $\mu \nabla^2 \vec{v}$: 점성에 의한 힘
  \item $\vec{f}$: 외부 힘(중력 등)
\end{itemize}

\subsection{역사}

\textbf{클로드 루이 나비에}(1785--1836)가 1822년에 처음 유도. \textbf{조지 가브리엘 스토크스}(1819--1903)가 1845년에 응력텐서를 이용한 명료한 재유도. 두 사람의 이름이 모두 붙은 이유다.

\subsection{밀레니엄 난제}

2000년 클레이 수학연구소가 7대 밀레니엄 난제 중 하나로 선정했다.

\begin{quotebox}
3차원 공간에서, 매끄러운 초기 조건이 주어졌을 때, 나비에-스토크스 방정식의 해가 항상 매끄럽게 영원히 존재하는가? 아니면 유한한 시간 안에 해가 폭발(blow-up)할 수 있는가?
\end{quotebox}

상금은 100만 달러. 비선형 항 $(\vec{v} \cdot \nabla)\vec{v}$ 때문에 중첩 원리가 성립하지 않으며, 이게 \textbf{난류(turbulence)}의 본질이다. 리처드 파인만은 난류를 ``고전물리학에 남은 가장 중요한 미해결 문제''라 불렀다.

% ====================================================================
\section{심혈관 시뮬레이션}
% ====================================================================

\subsection{왜 어려운 문제인가}

혈관 안의 흐름은 나비에-스토크스 방정식의 가장 도전적인 응용 중 하나다.

\textbf{혈관의 복잡성}: 대동맥 직경 2--3 cm, 모세혈관 5--10 $\mu$m로 3000배 이상의 스케일 차이. 혈관벽은 \textbf{탄성}이 있어 유체-구조 상호작용(FSI)으로 풀어야 한다.

\textbf{혈액의 복잡성}: 적혈구·백혈구·혈소판이 떠 있는 \textbf{현탁액}이며, 변형률에 따라 점성이 달라지는 \textbf{비뉴턴 유체}.

\textbf{시간적 변화}: 심박동에 의한 \textbf{맥동성(pulsatile)} 흐름.

\subsection{시뮬레이션 절차}

\begin{enumerate}[leftmargin=*, itemsep=2pt]
  \item \textbf{영상화}: CT 혈관조영술(CTA) 또는 MRA로 3D 영상 획득
  \item \textbf{메쉬 생성}: 수백만\textasciitilde 수천만 셀의 3D 메쉬 구축
  \item \textbf{경계조건}: 입구는 측정된 속도 프로파일, 출구는 0D 회로 모델
  \item \textbf{수치해석}: 매 시간 단계마다 모든 셀에서 나비에-스토크스 방정식 해석
  \item \textbf{임상 해석}: 벽전단응력(WSS), 압력 분포, 진동 전단지수(OSI), 유선과 와류 등 분석
\end{enumerate}

\subsection{임상 응용 사례}

\subsubsection{FFRct --- 관상동맥 협착 평가}

HeartFlow사의 대표적 상용 사례. 전통적 침습 카테터 검사 대신, \textbf{CT 영상만으로} 나비에-스토크스 방정식을 풀어 FFR(Fractional Flow Reserve) 값을 추정한다. 미국 FDA 승인, 영국 NHS, 일본 보험 적용.

\subsubsection{뇌동맥류 파열 위험 평가}

크기와 모양만으로 부족하므로, 동맥류 내부의 혈류 시뮬레이션으로 다음을 평가한다.
\begin{itemize}[leftmargin=*, itemsep=1pt]
  \item 낮고 진동하는 벽전단응력 (파열 위험 증가)
  \item 동맥류 내 복잡한 와류 패턴
  \item 좁고 빠른 입구 제트
\end{itemize}

\subsubsection{기타 응용}

\begin{itemize}[leftmargin=*, itemsep=1pt]
  \item 대동맥 박리·동맥류 추적 관찰
  \item 선천성 심장병의 가상 수술 계획 (폰탄 수술 등)
  \item TAVR 인공판막 설계 및 환자별 맞춤
  \item 스텐트 설계 및 시술 계획
\end{itemize}

\subsection{한계}

\begin{itemize}[leftmargin=*, itemsep=1pt]
  \item 정밀 FSI는 환자 한 명당 수 시간\textasciitilde 수일 소요
  \item 경계조건의 측정 불확실성
  \item 혈관벽의 비등방성·점탄성 모델링
  \item 일부 병적 상황의 난류
\end{itemize}

% ====================================================================
\section{유한체적법 (FVM)}
% ====================================================================

\subsection{발상: 보존법칙의 직접 적용}

FVM의 핵심 발상:

\begin{quotebox}
\textbf{미분방정식을 풀지 말고, 각 작은 부피 안에서 보존법칙을 직접 적용하자.}
\end{quotebox}

공간을 작은 \textbf{제어체적(control volume)}으로 쪼개고, 각 셀마다 다음을 강제한다.
\[
\text{셀 안 양의 변화율} = \text{셀 면을 통해 들어오는 플럭스의 합}
\]

\subsection{수학적 형식}

질량 보존 $\nabla \cdot \vec{v} = 0$을 셀 부피 $V$에 대해 적분하면 가우스 정리에 의해
\[
\int_V \nabla \cdot \vec{v}\, dV = \oint_{\partial V} \vec{v} \cdot \vec{n}\, dS = 0
\]

이산화하면
\[
\sum_{f} \vec{v}_f \cdot \vec{n}_f \cdot A_f = 0
\]

각 면에서 빠져나가는 양의 합이 0이라는 것이다.

\subsection{FEM과의 비교}

\begin{table}[h]
\centering
\small
\begin{tabular}{p{4cm}p{5cm}p{5cm}}
\toprule
& \textbf{유한요소법 (FEM)} & \textbf{유한체적법 (FVM)} \\
\midrule
무엇을 근사하나 & 해 자체를 다항식의 합으로 표현 & 셀에서 들어오고 나가는 플럭스 \\
미지수 위치 & 절점(꼭짓점) & 셀 중심 \\
출발점 & 변분 원리, 약한 형식 & 보존법칙의 적분 형식 \\
보존성 & 자동 보장 안 됨 & 이산화 단계에서 자동 보장 \\
복잡한 형상 & 매우 강함 & 가능하지만 까다로움 \\
주된 응용 & 구조역학, 고체역학 & 유체역학, 열전달, 충격파 \\
\bottomrule
\end{tabular}
\caption{FEM과 FVM의 비교}
\end{table}

\subsection{왜 유체에는 FVM이 표준인가}

\begin{enumerate}[leftmargin=*, itemsep=1pt]
  \item \textbf{질량 보존}: 비압축성 유체에서 메쉬 수준 보존 자동 보장
  \item \textbf{불연속 처리}: 충격파, 급격한 압력 변화에 강건
  \item \textbf{디버깅 용이}: 회계 장부 비유의 명료함
\end{enumerate}

ANSYS Fluent, OpenFOAM, Star-CCM+ 등 주요 CFD 솔버가 모두 FVM 기반이다.

\subsection{역사}

\textbf{FEM}: 1940--50년대 항공우주공학에서 시작(보잉). 1960년대 수학적 이론 정립.

\textbf{FVM}: 1970년대 패탱카(Suhas Patankar)의 \emph{Numerical Heat Transfer and Fluid Flow}(1980)로 표준화. SIMPLE 알고리즘 등.

% ====================================================================
\section{물리 정보 신경망 (PINN)}
% ====================================================================

\subsection{발상의 전환}

기존 수치해석은 메쉬를 만들고 각 셀에서 미지수를 정의하지만, PINN은 다르다.

\begin{quotebox}
\textbf{해 $u(x, t)$ 자체를 신경망 $u_\theta(x, t)$로 직접 표현한다.}
\end{quotebox}

여기서 $\theta$는 신경망의 가중치다. 메쉬도 셀도 요소도 없다. 입력 $(x, t)$를 받으면 출력 $u$가 나오는 매끄러운 함수일 뿐이다. 푸리에가 사인·코사인을 기저로 함수를 펼친 것의 현대판이며, 기저가 신경망의 비선형 변환으로 바뀐 셈이다.

\subsection{핵심 트릭: 자동미분}

신경망은 미분 가능한 함수다. 입력에 대한 출력의 미분을 \textbf{컴퓨터가 정확히} 계산할 수 있다. 따라서 $\dfrac{\partial u_\theta}{\partial t}$나 $\dfrac{\partial^2 u_\theta}{\partial x^2}$ 같은 도함수를 수치 차분 같은 근사 없이 정확히 얻을 수 있다. 이게 PINN을 가능하게 만든 핵심 기술이다.

\subsection{손실 함수: 물리를 학습 목표로}

PINN은 \textbf{물리 법칙 자체를 손실 함수}로 만든다. 도메인 안의 점들에서 방정식 잔차를 측정:
\[
L_{\text{physics}} = \frac{1}{M} \sum_j \left| \frac{\partial^2 u_\theta}{\partial t^2}(x_j, t_j) - c^2 \frac{\partial^2 u_\theta}{\partial x^2}(x_j, t_j) \right|^2
\]

여기에 초기조건·경계조건 손실, 데이터 손실을 더해 전체 손실:
\[
L = \lambda_1 L_{\text{data}} + \lambda_2 L_{\text{physics}} + \lambda_3 L_{\text{IC}} + \lambda_4 L_{\text{BC}}
\]

경사하강법으로 이를 최소화하면, 신경망이 방정식과 경계조건을 동시에 만족하는 함수로 학습된다.

\subsection{창시자}

마이지(Maziar Raissi), 페르디카리스(Paris Perdikaris), 카르니아다키스(George Karniadakis)의 2017--2019년 일련의 논문, 특히 2019년 \emph{Journal of Computational Physics} 논문이 결정적이었다.

\subsection{강점}

\begin{itemize}[leftmargin=*, itemsep=2pt]
  \item \textbf{역문제(inverse problem)}: 관측 데이터로부터 미지 파라미터 추정에 강력
  \item \textbf{메쉬 프리}: 복잡한 도메인에서 메쉬 생성 불필요
  \item \textbf{차원의 저주에 강함}: 고차원 PDE에서 강점
  \item \textbf{연속 해}: 임의의 점에서 즉시 값과 도함수 계산
  \item \textbf{데이터와 물리의 자연스러운 결합}
\end{itemize}

\subsection{약점}

\begin{itemize}[leftmargin=*, itemsep=2pt]
  \item \textbf{훈련 어려움}: 손실 가중치 균형 문제
  \item \textbf{정확도}: 잘 정립된 문제에서는 FEM/FVM이 보통 더 정확
  \item \textbf{충격파·난류}: 매끄러움 가정 때문에 불연속에 약함
  \item \textbf{재훈련 필요}: 새 문제마다 처음부터 훈련
  \item \textbf{이론적 보장}: 수렴성·안정성·오차 한계 이론이 아직 미완성
\end{itemize}

\subsection{FVM vs PINN: 연구 현황}

\subsubsection{FVM --- 성숙기}

이론은 1970--90년대에 정립. 현재는 \textbf{점진적 개선}이 중심:
\begin{itemize}[leftmargin=*, itemsep=1pt]
  \item 적응형 메쉬 미세화(AMR)
  \item 고차 정확도 스킴
  \item GPU 병렬화, 엑사스케일 컴퓨팅
  \item 불연속 갈러킨법 같은 하이브리드
\end{itemize}

ANSYS Fluent, OpenFOAM 등 산업 표준 솔버 정착. 연간 수십억 달러 규모의 CFD 시장.

\subsubsection{PINN --- 폭발기}

2019년 원조 논문 이후 매년 거의 두 배씩 논문 증가. 활발한 방향:

\begin{itemize}[leftmargin=*, itemsep=1pt]
  \item \textbf{변형들}: XPINN(도메인 분할), cPINN(보존형), fPINN(분수 미분), gPINN(그래디언트 강화), VPINN(변분형) 등
  \item \textbf{신경 연산자}: DeepONet, Fourier Neural Operator(FNO) --- PDE의 해 연산자 자체를 학습해 재훈련 문제를 해결
  \item \textbf{기초공업체}: DeepXDE, NVIDIA Modulus(이전 SimNet)
  \item \textbf{응용 확장}: 의학, 기후, 재료, 양자화학, 금융
\end{itemize}

\subsubsection{솔직한 비교}

\begin{notebox}
\textbf{연구 활동의 양과 신규성}: PINN/신경 연산자가 압도적으로 활발.

\textbf{산업 활용도}: 여전히 FVM이 압도적. 비행기 설계, 자동차 공력, 환자 진단 등 실용 도구는 거의 모두 FVM/FEM 기반.

\textbf{관계}: 대결이 아니라 \textbf{상호 보완}. 신경망을 FVM의 가속기로 사용하거나, FVM이 정답을 만들고 PINN이 학습해 추론을 단축시키는 방향.
\end{notebox}

% ====================================================================
\section{맺음말}
% ====================================================================

달랑베르가 1747년에 던진 진동하는 줄 한 가닥은, 270년이 지난 지금까지도 인류 수학과 공학의 흐름을 결정짓고 있다.

이 여정의 마디들을 다시 짚어보면:

\begin{itemize}[leftmargin=*, itemsep=2pt]
  \item 1747년 --- 달랑베르가 파동방정식과 그 해 $u = f(x-ct) + g(x+ct)$ 제시
  \item 1822년 --- 푸리에가 ``임의의 함수는 사인·코사인의 합'' 선언
  \item 1822/1845년 --- 나비에·스토크스가 유체 방정식 정립
  \item 1904년 --- 프란틀의 경계층 이론이 달랑베르의 역설 해결
  \item 1970년대 --- 패탱카의 FVM이 CFD의 표준으로
  \item 2000년 --- 클레이 연구소가 나비에-스토크스를 100만 달러 난제로
  \item 2019년 --- 카르니아다키스 등이 PINN 프레임워크 정립
  \item 현재 --- 신경 연산자, 머신러닝과 전통 수치해석의 융합
\end{itemize}

라그랑주의 다항식 기저, 푸리에의 사인·코사인 기저를 지나, 21세기는 \textbf{신경망이라는 새로운 기저}로 함수를 표현하는 시대로 들어선 셈이다. 본질은 250년 전 달랑베르가 던진 질문과 같다.

\begin{quotebox}
\centering
\textbf{매끄럽지 않은 자연 현상을, 어떤 매끄러운 도구들의 조합으로 표현할 것인가?}
\end{quotebox}

성당 계단에 버려진 갓난아기에서 시작된 이야기가, 오늘 환자의 심장 안 혈류를 보는 도구와 인공지능 시대의 새로운 수치해석으로 이어지고 있다. 진동하는 줄 한 가닥의 긴 메아리이다.

\vfill
\begin{center}
\rule{0.5\textwidth}{0.4pt}\\
\vspace{0.3cm}
\textit{Crowdy \& Claude 대화록, 2026}\\
\vspace{0.3cm}
\rule{0.5\textwidth}{0.4pt}
\end{center}

\end{document}
